\begin{abstract}
Prompt-and-validate pipelines---an ordinary chat endpoint, a Pydantic
model and a repair loop---are a common way to obtain structured output
from hosted language models, and their evaluation usually reduces to a
validity rate. We define a taxonomy of schema-conformance failures for
this paradigm, assigned by a deterministic classifier over the raw
completion and the validator's error, and apply it to
\pn{reported.attempts} generation attempts across \pn{reported.cells}
configurations spanning \pn{reported.models} models from two vendors,
one baseline schema with two targeted variants, and two task domains.
Three findings follow that a rate cannot express. First, models with
similar failure rates fail in different ways, and the dominant category
depends on the constraint surface and the task domain as well as the
model: relaxing length caps moves one model from a length-dominant to a
pattern-dominant profile, and revalidating its stored completions shows
that relaxing the caps is sufficient to eliminate those observed
violations. Second, the most frequent category, markdown wrapping,
conceals a further violation in \pn{pooled.reported.fence_compound_share}
of wrapped failures, from \pn{g.claude-sonnet-5@medium.fence_compound_share}
for one wrapping model to \pn{g.claude-haiku-4-5@medium.fence_compound_share}
for another, so fence removal is a partial repair whose effectiveness
varies substantially across configurations. For one model every first
attempt was fenced and no attempt after the second was; every run that did not complete ended on a length violation already present, beneath the fence, in the same field of its first attempt.
Third, because wrapping takes precedence in classification, prompt
variants that suppress fences raise the primary length-violation rate
without raising the combined rate of primary and concealed length
violations; primary labels alone would suggest a migration between
failure modes. We release the classifier, raw completions and annotation
record.
\end{abstract}
